\scanpage{72}
% archive sheet number: 71
% The entire opening paragraph is crossed out by large diagonal strokes.
\deleted{Nous allons examiner de plus près le cas où $\alpha,\beta$ sont des racines de polynômes $\frac12$ cyclotomiques $F_p(\alpha)=0,F_{\unsure{q|s}}(\beta)=0$ ($p,q\ge1$, sans exclure les cas $1,2$ qui correspondent \inserted{resp.} aux valeurs $2,-2$ pour $\alpha,\beta\ldots$), \deleted{\unsure{dans}} \inserted{et} le cas \uline{réel} $k=\mathbb{R}$. Si on veut obtenir \unsure{une|un} \deleted{\unsure{homéomorphisme|homéomorphie}} \inserted{\unsure{plongement}} \deleted{\unsure{image}} de la réalisation topologique de la carte combinatoire $C_{p,q}$, il faut prendre (pour $p$ resp. $q$ finis $\ge3$) des valeurs de la $\alpha$ resp. $\beta$ de la forme $\unsure{\zeta|\eta}+\zeta^{-1}$, où $\zeta=\exp 2i\pi/n$ ($n=p$ resp. $q$); d'autre part, $\alpha=-2\inserted{\text{ ou }\beta=-2}$ donne des croisements, donc si $p=\infty$ il faut $\alpha=2$, et $q=\infty$ il faut $\beta=2$.}

\[
\beta=2\qquad\alpha\ne-2
\]
\[
\gamma=2(3+\alpha)+2\alpha=2(2\alpha+3)
\qquad\gamma+2=2(2\alpha+4)=4(\alpha+2)
\]
\[
\begin{aligned}
\delta'&=4-4(\gamma+2)-\gamma^2=-(\gamma^2+4\gamma+4)=-(\gamma+2)^2\\
&=-16(\alpha+2)^2
\end{aligned}
\]
qui n'est nul que si $\alpha=-2$, cas que nous avons exclu\ldots
\[
\begin{aligned}
f_0(x,y,z)&=x^2+y^2+z^2-2(2\alpha+3)yz-2zx-2xy\\
&=\deleted{\unsure{xy|uxy}}(-x+y+z)^2-2(2\alpha+4)yz
\end{aligned}
\]
\[
f_0(x,y,z)=(-x+y+z)^2-4(\alpha+2)yz
\]
